\documentclass[12pt]{article}
\usepackage[margin=0.75in,top=0.8in,bottom=0.65in]{geometry}
\usepackage{lmodern}
\usepackage{microtype}
\usepackage{fancyhdr}
\usepackage{titlesec}
\usepackage{enumitem}
\usepackage{lastpage}
\usepackage[hidelinks]{hyperref}

\setlength{\parindent}{0pt}
\setlength{\parskip}{0.38em}
\setlength{\headheight}{26pt}
\titleformat{\section}{\large\bfseries\scshape}{}{0pt}{}
\titleformat{\subsection}{\normalsize\bfseries}{}{0pt}{}
\pagestyle{fancy}
\fancyhf{}
\fancyhead[C]{\footnotesize Lit Example Reader \quad Assignment ID: U1L5LE \quad Page \thepage\ of \pageref{LastPage}}
\renewcommand{\headrulewidth}{0.5pt}

\newenvironment{playpassage}{\begin{quote}\small\setlength{\parskip}{0.25em}}{\end{quote}}
\newcommand{\speaker}[1]{\textbf{#1.}}

\begin{document}

\begin{center}
{\LARGE\bfseries Week 5 Lit Example Reader}\par\vspace{0.05em}
{\large\scshape Julius Caesar: Brutus, Antony, and Formal Shape}\par\vspace{0.15em}
\rule{0.78\textwidth}{0.5pt}
\end{center}

\textbf{Purpose:} This reader gives Julius Caesar passages for applying Week 5's logic habit: identify mood, test formal validity, and separate a clean form from a proved premise.

\textbf{What to watch for:} Brutus and Antony both speak about honor, ambition, freedom, and death. The emotional language matters, but Week 5 asks first whether the form of the argument can carry the conclusion if the premises are granted.

\section*{Before the Passages: The Week 5 Logic Tool}

A categorical syllogism has three claims. Its \textbf{mood} names the form of the major premise, minor premise, and conclusion: A, E, I, or O. A beginner valid form is AAA in the first figure: All M are P; All S are M; therefore All S are P. Week 5's question is not yet whether every premise is true. First ask: \textit{if the premises were true, would the conclusion have to follow?}

\section*{Passage 1: Act 3, Scene 2 --- Brutus Explains the Killing}

\textbf{Context:} Brutus has joined the assassination of Caesar. He now faces the Roman citizens and tries to justify what the conspirators have done.

\begin{playpassage}
\speaker{Brutus} Romans, countrymen, and lovers! hear me for my cause,\par
and be silent, that you may hear: believe me for mine honour,\par
and have respect to mine honour, that you may believe.\par
If then that friend demand why Brutus rose against Caesar,\par
this is my answer:---Not that I loved Caesar less, but that I\par
loved Rome more. Had you rather Caesar were living and die all\par
slaves, than that Caesar were dead, to live all free men? As\par
Caesar loved me, I weep for him; as he was fortunate, I rejoice\par
at it; as he was valiant, I honour him: but, as he was ambitious,\par
I slew him. There is tears for his love; joy for his fortune;\par
honour for his valour; and death for his ambition.\par
Who is here so base that would be a bondman? If any, speak;\par
for him have I offended. Who is here so rude that would not be\par
a Roman? If any, speak; for him have I offended.
\end{playpassage}

\subsection*{Reader Notes}

Here is your most assisted example. One possible formal map is: All men whose ambition would enslave Rome deserve death for Rome's freedom; Caesar was a man whose ambition would enslave Rome; therefore Caesar deserved death for Rome's freedom. That map has an AAA shape: All M are P; All S are M; therefore All S are P. The form is strong if the terms are held steady.

\textbf{Reading task:} Underline Brutus's conclusion. Circle the word \textit{ambitious}. Mark the sentence that makes the crowd choose between slavery and freedom.

\newpage

\section*{Passage 2: Act 3, Scene 2 --- Antony Presses the Middle Term}

\textbf{Context:} Antony speaks after Brutus. He does not begin by diagramming a syllogism. Instead, he repeats Brutus's key term and gives evidence that pressures it.

\begin{playpassage}
\speaker{Antony} The noble Brutus\par
Hath told you Caesar was ambitious:\par
If it were so, it was a grievous fault,\par
And grievously hath Caesar answer'd it.\par
Here, under leave of Brutus and the rest---\par
For Brutus is an honourable man;\par
So are they all, all honourable men---\par
Come I to speak in Caesar's funeral.\par
He was my friend, faithful and just to me:\par
But Brutus says he was ambitious;\par
And Brutus is an honourable man.\par
He hath brought many captives home to Rome\par
Whose ransoms did the general coffers fill:\par
Did this in Caesar seem ambitious?\par
When that the poor have cried, Caesar hath wept:\par
Ambition should be made of sterner stuff:\par
Yet Brutus says he was ambitious;\par
And Brutus is an honourable man.\par
You all did see that on the Lupercal\par
I thrice presented him a kingly crown,\par
Which he did thrice refuse: was this ambition?\par
Yet Brutus says he was ambitious.
\end{playpassage}

\subsection*{Reader Notes}

Antony may not give a cleaner syllogism than Brutus. He gives counter-evidence. Watch whether his reasoning proves a universal negative, proves a weaker particular claim, or simply makes Brutus's minor premise doubtful. Decide how much his evidence can formally support.

\textbf{Reading task:} Number Antony's three pieces of evidence. Beside each, write whether it attacks the term \textit{ambitious}, the conclusion that Caesar deserved death, or both.

\section*{How to Use These Passages on the Worksheet}

The worksheet asks you to apply the form test. Do not skip straight to whether you like Brutus or Antony. First identify the shape. Then decide what the form proves and what still depends on evidence.

\end{document}
